\section{The even transfer sector}\label{sec:parity}

The trace counts every excitation, whereas the two-point function sees only
states created by the spin field. We first show that a lowest transfer
excitation is odd under a component reflection, with a quantitative overlap
controlled by its eigenvalue. The main ingredient is a finite-graph
correlation bound: the covariance of two even spin monomials is quadratic in
the largest two-point correlation between their supports.

For this section, consider any finite graph with spins $q_x\in S^2$ and law
\begin{equation}\label{par:graph-law}
 \frac1Z\exp\!\left(\sum_{xy}J_{xy}q_x\cdot q_y\right)
       \prod_x d\sigma(q_x),\qquad 0\le J_{xy}<\infty,
\end{equation}
where $\sigma$ is sphere probability measure. A list of sites records one
entry per factor of a monomial; repeated sites remain separate entries.

\begin{proposition}[Quadratic covariance bound]\label{par:covariance}
Let $P$ and $Q$ be monomials in spin components for the law
\eqref{par:graph-law}, of even total degrees $r$ and $s$, with nonempty
site lists $\mathcal A$ and $\mathcal B$. Put
\[
 G=\max_{x\in\mathcal A,\ y\in\mathcal B}\E(q_x^3q_y^3).
\]
There is a constant $C_{r,s}$, depending only on the two degrees, such that
\begin{equation}\label{par:quadratic-bound}
             |\Cov(P,Q)|\le C_{r,s}G^2.
\end{equation}
If either monomial is constant, its covariance is zero.
\end{proposition}

We prove the proposition in three stages. A component decomposition first
gives correlation signs. A positive pairing expansion then proves the bound
for single-component monomials. Finally, a Fourier comparison reduces mixed
components to that case. These arguments use the correlation-inequality
methods of Griffiths, Kelly--Sherman, and Ginibre
\cite{Griffiths1967,KellySherman1968,Ginibre1970}; the component comparisons
are also related to the multicomponent rotator inequalities of
Dunlop~\cite{Dunlop1976}.

\subsection{Component decomposition and correlation signs}

Write, independently at each site,
\begin{equation}\label{par:angles}
 q_x=(\cos\alpha_x\cos\vartheta_x,
       \cos\alpha_x\sin\vartheta_x,
       \sin\alpha_x u_x),
 \quad 0\le\alpha_x\le\frac\pi2,
 \quad u_x\in\{-1,1\},\quad \vartheta_x\in\R/2\pi\Z.
\end{equation}
The one-site reference law is the product of $\cos\alpha_x\,d\alpha_x$
and the uniform laws of $u_x$ and $\vartheta_x$. Given $\alpha$, the signs
and angles are independent, with respective Ising and plane-rotator
couplings
\begin{equation}\label{par:conditional-couplings}
 J^{\mathrm I}_{xy}=J_{xy}\sin\alpha_x\sin\alpha_y,
 \qquad
 J^{\mathrm R}_{xy}=J_{xy}\cos\alpha_x\cos\alpha_y.
\end{equation}
We order angle arrays coordinatewise. A law is \emph{positively associated}
if the covariance of every pair of bounded increasing functions is
nonnegative.

We recall explicitly the plane-rotator inequality needed below. For any
integer vectors $k,l$ indexed by the graph vertices,
\begin{equation}\label{par:rotator-inequalities}
 \E\cos(k\cdot\vartheta)\ge0,
 \qquad
 \Cov\bigl(\cos(k\cdot\vartheta),
            \cos(l\cdot\vartheta)\bigr)\ge0.
\end{equation}
The first inequality follows by expanding each bond exponential into its
nonnegative Fourier coefficients. Here is the replica proof of the second,
in the form of Ginibre's method. With independent copies $\vartheta,
\vartheta'$, twice the covariance is the expectation of the product of the
two observable differences. The substitution
\[
             \vartheta=U+V,\qquad \vartheta'=U-V
\]
preserves product Haar integration: it is a surjective homomorphism of the
product torus. Each difference is
$-2\sin(k\cdot U)\sin(k\cdot V)$, while the replicated bond exponent is
$2J_{xy}\cos(U_x-U_y)\cos(V_x-V_y)$. Expanding its exponential gives
nonnegative coefficients times squares of real integrals, proving
\eqref{par:rotator-inequalities}. In particular every cosine mean is
nondecreasing in each nonnegative coupling, by differentiation.
The zero-field Ising GKS inequalities give the corresponding nonnegative
means and covariances for sign monomials
\cite{Griffiths1967,KellySherman1968}.

\begin{lemma}[Association and single-component signs]\label{par:association}
The marginal law of $\alpha$ in \eqref{par:angles} is positively associated.
For even-degree monomials each using a single component, covariances are
nonnegative when the components agree and nonpositive when they differ.
The same covariance signs hold for the two-component plane rotator.
\end{lemma}

\begin{proof}
Relative to the product angle reference law, the density of $\alpha$ is
proportional to $Z_{\mathrm I}(\alpha)Z_{\mathrm R}(\alpha)$. For either
conditional model, write $O_e$ for its bond observable and $J_e(\alpha)$
for its coupling. At distinct vertices $x,y$,
\begin{equation}\label{par:mixed-partials}
 \partial_x\partial_y\log Z
 =\sum_e\E(O_e)\,\partial_x\partial_yJ_e
  +\sum_{e,f}\Cov(O_e,O_f)\,
                      \partial_xJ_e\,\partial_yJ_f\ge0.
\end{equation}
Indeed all first derivatives of the sine couplings are nonnegative and all
first derivatives of the cosine couplings are nonpositive. A coupling's
only nonzero mixed derivative at distinct sites is at its endpoints, where
it is nonnegative in both cases. All bond means and covariances in
\eqref{par:mixed-partials} are nonnegative by the inequalities just proved
or recalled. Thus the density satisfies the FKG lattice condition, which
implies association~\cite{FKG1971}. One may restrict the angles to compact
subintervals of $(0,\pi/2)$, discretize, and pass to the limit; this also
handles the endpoints of the angle interval.

To obtain the same fact for the plane rotator, write its two coordinates
as $(\cos\gamma\,\varepsilon,\sin\gamma\,\varepsilon')$, where
$0\le\gamma\le\pi/2$ and the two signs are conditionally independent
ferromagnetic Ising systems. The calculation
\eqref{par:mixed-partials} proves association of $\gamma$. The conditional
mean of a single-component monomial is nonnegative. Including its
amplitude factors, it decreases with $\gamma$ for the first coordinate
and increases for the second. Conditional independence and association
therefore give nonpositive covariance between unlike components. For like
components, the conditional covariance is nonnegative by GKS, and the
covariance of conditional means is nonnegative by association. This proves
both plane-rotator assertions.

In the three-component decomposition, the conditional mean of a
color-$3$ monomial is likewise nonnegative and increasing in $\alpha$.
The conditional mean of a color-$1$ monomial is nonnegative and decreasing:
a product of cosines is a nonnegative linear combination of
$\cos(k\cdot\vartheta)$, and both the amplitude factors and the rotator
couplings decrease with $\alpha$. Conditional independence now gives the
nonpositive covariance between colors $3$ and $1$. For like colors, use
GKS or \eqref{par:rotator-inequalities} for the conditional covariance and
association for the covariance of conditional means. Component symmetry
gives every remaining choice of colors.
\end{proof}

\subsection{Pairings and the reduction of mixed components}

\begin{proof}[Proof of Proposition~\ref{par:covariance}]
Expand each bond exponential in \eqref{par:graph-law} into powers of the
dot product. The uniform sphere moment tensor of order $2h$ is
\[
 \int_{S^2}q^{i_1}\cdots q^{i_{2h}}\,d\sigma(q)
 =\frac1{3\cdot5\cdots(2h+1)}
       \sum_{\pi}\prod_{\{b,c\}\in\pi}\one_{i_b=i_c},
\]
where the sum is over pairings of the $2h$ labeled factors; odd tensors
vanish. Following the contracted bond lines through these site pairings
pairs the external factors, with closed lines contributing positive color
sums. Consequently, for an arbitrary color assignment to the combined
lists $\mathcal A,\mathcal B$, the joint moment has the form
\begin{equation}\label{par:pairings}
 \E\prod_{b=1}^{r+s}q_{x_b}^{i_b}
   =\sum_{\pi}d_\pi
       \prod_{\{b,c\}\in\pi}\one_{i_b=i_c},
 \qquad d_\pi\ge0.
\end{equation}
The coefficients do not depend on the external colors. The finite graph
and finite couplings ensure absolute convergence of the exponential
expansion; all compatible contracted terms are nonnegative. This
justifies collecting them by their external pairing, including when sites
are repeated.

\emph{Crossing pairs.}
Call a pair crossing if it joins the two lists. Since both list sizes are
even, a pairing with a crossing pair has at least two. For any such
pairing we claim
\begin{equation}\label{par:pair-coefficient}
                         d_\pi\le G^2.
\end{equation}
Assign color $3$ to the endpoints of one crossing pair, color $2$ to the
endpoints of another, and color $1$ to all remaining endpoints. The
resulting moment contains $d_\pi$ as a nonnegative summand. Condition on
$\alpha$. In the rotator model, Lemma~\ref{par:association} bounds the
moment of the color-$2$ pair times the remaining color-$1$ monomial by
the product of their means, and hence by the color-$2$ pair mean, because
the color-$1$ monomial's mean lies in $[0,1]$. Include the factors
 $\cos\alpha_x$ in this statement. The resulting pair mean is nonnegative
and decreasing in $\alpha$. Indeed, rotator invariance gives
\[
 \E(\sin\vartheta_x\sin\vartheta_y\mid\alpha)
       =\tfrac12\E\bigl(\cos(\vartheta_x-\vartheta_y)\mid\alpha\bigr),
\]
also when $x=y$, and both this cosine mean and the amplitude factors
decrease with $\alpha$. The color-$3$ pair mean, including its
sine factors, is nonnegative and increasing. Association bounds the
average of their product above by the product of their averages. Each
average is a two-point function between the lists, at most $G$ by component
symmetry. This proves \eqref{par:pair-coefficient}.

Consider next single-component monomials on the two lists. Compare the
case where their colors agree with the case where they differ. Their
products of means agree by symmetry, while their joint moments differ
by exactly the sum of the crossing contributions in
\eqref{par:pairings}. The two covariances have opposite weak signs by
Lemma~\ref{par:association}. Thus each absolute covariance is at most
the number of pairings times $G^2$. The same conclusion holds for any
even sublists of $\mathcal A$ and $\mathcal B$, with the same $G$.
If either $P$ or $Q$ has odd degree in at least one color, its own mean
vanishes. Every compatible pairing for $PQ$ then crosses the two lists,
so \eqref{par:pair-coefficient} also proves the required bound in this case.

It remains to treat monomials even in every color. This is the only case
where subtracting the product of mixed-component means cannot be read
directly from the pairing expansion. Write $P=W_PV_P$, where $W_P$ uses
color $3$ and $V_P$ uses colors $1,2$. Let $V_P^*$ replace every color in
$V_P$ by color $1$, and define
\[
 a_P=\E(W_P\mid\alpha),\qquad
 b_P=\E(V_P\mid\alpha),\qquad
 B_P=\E(V_P^*\mid\alpha).
\]
Use the corresponding notation for $Q$. We have $0\le a_P,B_P\le1$
and $|b_P|\le1$; $a_P$ is increasing and $B_P$ decreasing.
Conditional independence of signs and rotator angles gives
$\E(P\mid\alpha)=a_Pb_P$. Thus the total covariance identity reads
\begin{equation}\label{par:total-covariance}
 \Cov(P,Q)=\E\Cov(P,Q\mid\alpha)
                    +\Cov(a_Pb_P,a_Qb_Q).
\end{equation}
We bound the conditional fluctuations first, then the covariance of the
conditional means. A Fourier comparison with the single-color monomials
$V_P^*,V_Q^*$ will control both terms.

\emph{The common Fourier comparison.}
Factor out the common amplitude $A_P(\alpha)$, a product of cosines,
and expand the coordinate factors:
\begin{equation}\label{par:fourier-domination}
 V_P=A_P(\alpha)\sum_k c_k\cos(k\cdot\vartheta),
 \qquad
 V_P^*=A_P(\alpha)\sum_k d_k\cos(k\cdot\vartheta),
 \qquad |c_k|\le d_k.
\end{equation}
Each factor has two exponential terms with coefficients of absolute
value $1/2$. Replacing sines by cosines replaces their phases by positive
coefficients. Since there are evenly many sine factors, the coefficients
are real and invariant under $k\mapsto-k$. Grouping terms at a common
frequency preserves the displayed domination, even at repeated sites.

\emph{Conditional fluctuations.}
Conditional independence gives
\begin{align*}
 \Cov(P,Q\mid\alpha)
  ={}&\Cov(W_P,W_Q\mid\alpha)\,
                             \E(V_PV_Q\mid\alpha)\\
    &+a_Pa_Q\Cov(V_P,V_Q\mid\alpha).
\end{align*}
The first conditional covariance is nonnegative. The Fourier comparison
\eqref{par:fourier-domination} and the entrywise nonnegative cosine
covariances in \eqref{par:rotator-inequalities} imply
\[
 |\Cov(V_P,V_Q\mid\alpha)|
               \le\Cov(V_P^*,V_Q^*\mid\alpha).
\]
All remaining factors have absolute value at most one. Averaging thus
bounds $|\E\Cov(P,Q\mid\alpha)|$ by
\[
 \E\Cov(W_P,W_Q\mid\alpha)
       +\E\Cov(V_P^*,V_Q^*\mid\alpha).
\]
Each summand is at most its full single-component covariance: the
omitted covariance of conditional means is nonnegative, since both
$a$'s increase and both $B$'s decrease. The single-component bounds
therefore control the first term of \eqref{par:total-covariance}
by $C_{r,s}G^2$.

\emph{Covariance of conditional means.}
It remains to control $\Cov(a_Pb_P,a_Qb_Q)$. The functions $b_P,b_Q$
need not be monotone, but their changes are bounded by those of
$B_P,B_Q$. Indeed, \eqref{par:fourier-domination} shows that
$B_P+b_P$ and $B_P-b_P$ are nonnegative combinations of cosine means,
multiplied by $A_P$. Both functions are decreasing, by
\eqref{par:rotator-inequalities} and the decreasing rotator couplings.
For ordered angle arrays $\alpha\le\widetilde\alpha$, it follows that
\[
 |b_P(\widetilde\alpha)-b_P(\alpha)|
       \le B_P(\alpha)-B_P(\widetilde\alpha).
\]
Set $Y_P=a_Pb_P$ and $D_P=a_P-B_P$. Combining this estimate with the
increase of $a_P$ and the bounds $a_P,|b_P|\le1$ gives
\begin{align*}
 |Y_P(\widetilde\alpha)-Y_P(\alpha)|
 &\le a_P(\widetilde\alpha)-a_P(\alpha)
                  +B_P(\alpha)-B_P(\widetilde\alpha)\\
 &=D_P(\widetilde\alpha)-D_P(\alpha).
\end{align*}
In particular $D_P+Y_P$ and $D_P-Y_P$ are increasing.
Apply association to these functions and $D_Q\pm Y_Q$. Adding the two
inequalities with matching signs, and then the two with opposite signs,
gives
\begin{equation}\label{par:conditional-mean-bound}
               |\Cov(Y_P,Y_Q)|\le\Cov(D_P,D_Q).
\end{equation}
Every term in the expansion
\[
 \Cov(D_P,D_Q)
 =\Cov(a_P,a_Q)+\Cov(B_P,B_Q)
       -\Cov(a_P,B_Q)-\Cov(B_P,a_Q)
\]
is controlled by a single-component covariance already bounded above.
For the two like-color terms, conditional covariances are nonnegative,
so the covariances of conditional means are at most the full covariances.
For the unlike-color terms, conditional independence identifies the
covariances of conditional means with the full covariances exactly.
All their lists are even sublists of the original two lists. Hence
\eqref{par:conditional-mean-bound} bounds the second term of
\eqref{par:total-covariance} by $C_{r,s}G^2$ as well. Combining the two
bounds proves \eqref{par:quadratic-bound}.
\end{proof}

\subsection{Spectral consequence and overlap}

Return to the transfer operator at fixed cutoff $N$ and circumference $w$.
Write $T=T_{N,w}$, $\Omega=\Omega_{N,w}$, and
$\tau=e^{-ae_N(w)}\in(0,1)$. A row function is \emph{even} or \emph{odd}
according to its parity under simultaneous inversion of every row spin.
The vacuum $\Omega$ is even by its uniqueness and positivity.

\begin{proposition}[Even-sector bound and odd overlap]\label{par:even-gap}
On the even subspace perpendicular to the vacuum,
\begin{equation}\label{par:even-norm}
          \bigl\|T\big|_{\mathrm{even}\cap\Omega^\perp}\bigr\|
                         \le\tau^2.
\end{equation}
There is a real unit $\tau$-eigenfunction $\psi$ that is odd under
simultaneous spin inversion and, after a component permutation, odd under
the reflection of component $3$ alone. It satisfies
\begin{equation}\label{par:overlap}
                    \inner{\Omega}{|\psi|}^2\ge\frac{\tau}{1+\tau}.
\end{equation}
\end{proposition}

\begin{proof}
The finite-graph covariance bound passes to the infinite-time cylinder
by transfer convergence. For sites on two rows separated by $j$ time
steps, their single-component correlation is at most $\tau^j$: the vectors
$q_x^i\Omega$ are centered, have norm at most one, and $T$ has norm $\tau$
on $\Omega^\perp$. Proposition~\ref{par:covariance} therefore gives, for
every real even row polynomial $P$ and its centered vector
$v_P=(P-\inner{\Omega}{P\Omega})\Omega$,
\begin{equation}\label{par:polynomial-decay}
                   0\le\inner{v_P}{T^jv_P}\le C_P\tau^{2j}.
\end{equation}
Here we have summed the finitely many monomial covariance bounds; the
constant need not be uniform in $P,N$, or $w$.

Polynomials in row components are uniformly dense in the continuous
functions on the compact row space. Symmetrizing gives density of even
polynomials in the even subspace. The continuous positive function
$\Omega$ has a positive minimum, so multiplication by $\Omega$ is bounded
and invertible. Hence the complex linear span of the vectors $v_P$ is dense in the even
vacuum complement. Since the spectral measure of each $v_P$ is positive,
\eqref{par:polynomial-decay} excludes any of its mass above $\tau^2$:
positive mass in $[\tau',1]$ with $\tau'>\tau^2$ would give a lower bound
proportional to $(\tau')^j$. Density proves \eqref{par:even-norm}.

Compactness of $T$ implies that $\tau$ is an eigenvalue. As $\tau>\tau^2$,
its eigenspace is odd. The three component reflections commute with each
other and with the real operator $T$, so one may choose a real unit
eigenfunction with a definite parity under each reflection. Their product
is simultaneous inversion, so at least one component parity is odd;
relabel that component as $3$.

The modulus $|\psi|$ is even. Put $b=\inner{\Omega}{|\psi|}^2$. Positivity
of the transfer kernel and \eqref{par:even-norm} give
\[
 \tau=\inner{\psi}{T\psi}
       \le\inner{|\psi|}{T|\psi|}
       \le b+(1-b)\tau^2.
\]
Rearranging yields $b\ge\tau/(1+\tau)$, as claimed.
\end{proof}

Thus the lowest excitation has a bounded, component-odd detector:
$F=\sgn\psi$, with $F=0$ on the zero set, satisfies $|F|\le1$ and
$\inner{\Omega}{F\psi}=\inner{\Omega}{|\psi|}$. The next section uses
this detector to turn the spectral overlap into a connection probability.
